\section*{The Propers: Detailed Commentary}

\subsection*{Strength, shepherding, and holy fear\enspace\elementref{Entrance, Collect}}

The Entrance adapts Psalm 28's closing confession: the majestic Lord is also strength, refuge, savior, and shepherd. The image protects courage from self-reliance. The people facing witness are an inheritance carried by another, not a heroic elite sustained by temperament.

The full psalm moves through storm imagery and repeated attribution of glory before reaching this pastoral conclusion. Augustine's \emph{Exposition on Psalm 28} receives that movement sacramentally and ecclesially: the mighty voice that shakes creation is not separated from the Lord who gives strength and peace to his people. Whatever its allegorical details, Augustine's broad sequence prevents a sentimental shepherd image. Divine care belongs to the same holy majesty before which human strength is relativized. The antiphon adapts the closing verses; it does not appoint Augustine's allegory or the complete psalm.

The Collect's paired fear and love rejects both servile terror and sentimental familiarity. Reverence recognizes God's holiness and judgment; love cleaves to his name and trusts his governance.

\subsection*{Jeremiah's lament and Matthew's disclosure\enspace\elementref{First Reading, Psalm, Gospel}}

Jeremiah hears denunciation from every side, even from those watching for his fall. Yet he submits his cause to the One who tests the just and sees inward truth. The reading ends in praise, but the surrounding chapter proceeds into curse and anguish. Scripture therefore permits the persecuted believer to confess trust without pretending that trauma produces simple emotional serenity.

Jerome's \emph{Commentary on Jeremiah}, book IV, comments directly on 20:9--13 in Latin. He treats the unspoken word as a fire that the prophet cannot contain, sees surrounding reproach as a trial familiar to teachers, and notices that former friends wait for a fall. At verse 11 he reads the Lord's presence as that of a strong warrior; at verse 12 he does not sanitize the prophet's wish to see vindication, but identifies the human frailty and impatience disclosed by it. At verse 13 he assigns the rescue to grace rather than personal merit. Jerome does not settle prophetic psychology, but his verse-by-verse attention guards three features easily lost in thematic summaries: compelled speech, betrayal by intimates, and praise that credits God rather than the sufferer's prowess.

Jerome's sequence also prevents two opposite appropriations. Because the word has become inward fire, Jeremiah cannot be reduced to a speaker who simply chose controversy; because Jerome names impatience in the desire for visible vengeance, prophetic vocation does not place every affect or inference beyond judgment. The Lord who tests heart and mind is the prophet's defender and examiner. This combination makes the reading usable for ecclesial self-critique: threatened witnesses may seek God's judgment, while communities must still test whether speech, motive, and desired remedy accord with the God invoked.

The selected verses begin with conspiracy rather than the preceding account of the word burning within Jeremiah, and they stop at praise before the curse of the prophet's birth. Both contexts bear on the reading. The liturgy proclaims a real confession of deliverance, but its boundary does not mark the end of Jeremiah's anguish. Jerome's sequence itself confirms how closely confidence and impatience stand together.

Psalm 69 answers in the first person. Reproach touches the servant because of zeal for God's house; the petition seeks an answer in divine kindness; the poor are finally summoned to praise. Augustine hears the suffering Christ and his body within a psalm whose literal voice is Israelite lament.

Augustine's direct exposition repeatedly passes between Christ the head and suffering members, especially where reproach for God's cause is concerned. That reception shows why Christian tradition can pray the lament christologically even though not every historical detail first described Jesus. It also resists a cult of solitary heroism: the voice of the afflicted is heard within a body, and the psalm's final summons lets the poor join praise. The responsorial use is officially correlated with Jeremiah.

Matthew's repeated prohibition of fear is reasoned. What is concealed will be disclosed; disciples cannot finally secure truth through silence. Human persecutors possess limited power. God knows creatures down to sparrows and hairs, and Christ will acknowledge those who acknowledge him. These claims establish divine judgment and care, not present immunity.

Chrysostom's \emph{Homily 34 on Matthew} distinguishes the reasons Christ gives rather than reducing them to repetition. Hidden truth will be manifested; enemies cannot kill the soul; providence attends creatures of very small market value; and Christ will confess his confessors before the Father. Courage therefore arises from truth about the limits of persecuting power and the scope of divine knowledge. It is not an assertion that bodies do not suffer. Chrysostom's further comment that providence does not mean God positively causes every sparrow's fall is especially important: divine care is not a vocabulary for calling violence good.

Cyprian's \emph{Testimonies} III.4 later gathers Matthew 10:28 and the confession sayings for communities formed under the possibility of persecution. This is martyrial reuse rather than continuous exposition. It shows that the passage became a grammar of public fidelity; that grammar does not settle every disputed speech or private danger. Confession concerns Christ and truth under grace; it does not forbid prudent withdrawal, legal protection, or safeguarding.

The official first-reading--Gospel correlation joins two texts that keep their own settings. Jeremiah's enemies watch for a misstep; Matthew's disciples are told that concealed truth will be disclosed and that Christ will acknowledge faithful confessors. These motifs genuinely converse in the Lectionary. They do not turn Jeremiah into a coded prediction of Matthew 10, nor do they make every institutional conflict a prophetic persecution. The texts themselves supply tests: Jeremiah submits his cause to the God who sees heart and mind, while Jesus orders fear by divine judgment. Neither authorizes a speaker to declare himself righteous merely because he is opposed.

Chrysostom's distinction is precise: nothing falls outside God's knowledge and providential ordering, yet the text does not require saying that God positively effects every sparrow's fall. Catholic doctrine likewise distinguishes divine first causality, creaturely causes, permission of evil, and God's ability to bring good without making evil good. That distinction is decisive whenever this Gospel is preached to vulnerable persons.

\subsection*{The Spirit's testimony and the disciple's speech\enspace\elementref{Acclamation}}

John 15 places witness amid hatred and expulsion. The Spirit of truth testifies to Christ; disciples also testify because they have been with him. The acclamation therefore protects Matthew 10 from individualism. Public confession is ecclesial and pneumatological before it is a test of private nerve. The verse welcomes Christ's Gospel word; it is not another semi-continuous reading.

Augustine's \emph{Tractate 92 on John} expounds John 15:26--27 directly. He does not set the Spirit's testimony against the apostles' speech: the Spirit bears witness and makes them capable of bearing witness. Peter's passage from fearful denial to public confession provides Augustine's concrete example. Augustine's reading sharpens the acclamation's grammar. “You also testify” is not autonomous courage added alongside divine action; apostolic testimony is enabled from within by the promised Spirit and grounded in having been with Jesus. Since the acclamation quotes only parts of the two verses, its liturgical function remains a concise welcome to the Gospel rather than a substitute reading of the farewell discourse.

\subsection*{The gift is not like the trespass\enspace\elementref{Second Reading}}

Romans 5:12--15 begins a dense comparison whose grammar continues beyond the appointed boundary. Paul traces sin and death through Adam and then interrupts any expectation of equal symmetry: Christ's gracious gift is not like the trespass. Chrysostom emphasizes this excess. One act's devastation becomes the occasion for displaying a gift whose abundance cannot be represented as a mere zero-sum cancellation.

Aquinas's \emph{Commentary on Romans} V, lectio 3, attends to the difficult opening comparison and to Paul's interruption before the completed Adam--Christ parallel. The entrance of sin and death names a solidarity no individual created alone, while the gift comes through the one man Jesus Christ in an abundance not measured by the offense. Augustine's \emph{Against Two Letters of the Pelagians} IV.8 joins Romans 5:10--12 in an anti-Pelagian argument: those needing reconciliation are not merely neutral creatures lacking instruction, and grace cannot be reduced to natural capacity. Both are direct later receptions with their own controversies. Neither licenses treating inherited sin as personal blame for another's deed, nor does either erase Paul's insistence on the superabundance of Christ's gift.

Verse 12's final clause has generated major grammatical and doctrinal debate, but the appointed reading itself immediately supplies a corporate horizon: sin enters through one, death spreads, and all sin. The appointed verses therefore exclude two reductions---an account of Adam that makes personal sins irrelevant, and an individualism that cannot name a wounded solidarity preceding personal choice. The reading's decisive word is nevertheless not an abstract anthropology but the contrasting gift in verse 15. Paul introduces Adam to magnify Christ's abundance, not to leave the assembly contemplating ruin as the last reality.

The appointed reading ends at verse 15, before Paul completes several contrasts. This boundary limits what the reading itself teaches. The congregation hears the universality of death and the decisive inequality between trespass and gift; it does not hear the whole subsequent account of judgment, justification, obedience, and reigning grace. Those verses clarify Paul's grammar but are not part of the proclamation.

This apostolic strand remains semi-continuous. It is no allegory in which Adam equals persecutor or Christ equals Jeremiah. Its liturgical contribution is plural: the same assembly hearing the demand for witness also hears that salvation originates in a gift exceeding inherited ruin.

\subsection*{Reconciliation, praise, and nourishment\enspace\elementref{Prayer over Offerings, Communion, Prayer after Communion}}

The Prayer over the Offerings identifies the sacrificial movement as reconciliation and praise while asking that those offering be purified. Worship is therefore not a bribe paid by the endangered, nor a reward for courage; it is God's reconciling gift received and returned in praise.

Communion option A places every creature before the giver of food. Option B names the Good Shepherd whose self-gift distinguishes him from the hireling who abandons the flock. The Prayer after Communion then asks renewal through the precious Body and Blood toward redemption. The alternatives thus emphasize either creaturely dependence or the Shepherd's sacrificial care while converging in Christ's nourishing gift.

Augustine's \emph{Exposition on Psalm 145} receives the first branch's praise as a schooling away from trust in mortal rulers and toward the Creator who feeds, frees, raises, and keeps faith. The antiphon selects the eyes of all and God's timely gift of food; its psalmic context locates nourishment within a larger pattern of justice and care. That context prevents the verse from becoming a prosperity promise and prevents Eucharistic use from erasing ordinary creaturely dependence.

The verse's universal gaze also resists an exclusively inward Eucharistic reading: creatures await sustenance from their maker, and sacramental nourishment trains rather than abolishes attention to hungry bodies. Augustine's full-psalm reception places food among God's faithful acts toward the oppressed and bowed down. The antiphon remains one branch, but its context makes gratitude socially answerable.

For option B, Augustine's \emph{Tractate 46 on John} directly distinguishes the Good Shepherd, who gives his life, from the hireling, who flees when the wolf comes. \emph{Tractate 47} then treats the exact continuation in John 10:14--15. Christ knows his own and knows the Father; other shepherds act only as his members and enter through him, while his self-offering establishes their pattern without making them alternative doors. The appointed antiphon joins knowledge and self-gift, so pastoral familiarity without costly care is inadequate, as is sacrifice without truthful knowledge of the flock.

Cyril of Jerusalem's \emph{Catechetical Lecture 10}, section 5, receives the same saying christologically: Christ is sheep according to the humanity in which he is offered and Shepherd according to divine loving-kindness. Gregory Nazianzen's \emph{Oration 18} uses the shepherd's knowing, defending, and guiding of the flock as a pastoral measure. Augustine supplies direct exposition; Cyril and Gregory supply illuminating reuse. Together they make Communion B more than a generic picture of kindness: the Shepherd's care is self-giving, discerning, and answerable for the flock. Which alternative was enacted is unresolved.

\section*{The Propers: Notable and Quotable}

{\small
\begin{itemize}
\item \textbf{“Not one of them falls” --- Hamlet walks into a fatal duel.\enspace\elementref{Gospel}} In \emph{Hamlet} 5.2, the prince dismisses augury because “there is a special providence in the fall of a sparrow,” then accepts the fencing match that kills him. Shakespeare preserves Matthew's tiny bird while moving the saying into tragic dramatic irony: confidence in providence does not mean escape from death.

\Needspace{4\baselineskip}
\item \textbf{“All the hairs of your head are counted” --- and a bald man's back-order.\enspace\elementref{Gospel}} A 1900 college magazine reprinted “The Back Numbers”: a bald worshipper hears that every hair is numbered and prays, “O Lord, send me down the back numbers.” The joke makes divine enumeration sound like missing installments of a periodical.

\item \textbf{“The hairs of your head are counted” --- a quantified-self performance.\enspace\elementref{Gospel}} The 2017--2018 media-arts work \emph{The Hairs of Your Head Are Numbered} measured bodies, turned audience data into live visualization, and questioned tracking and manipulation. The exact Gospel-like title becomes an uneasy cultural echo for bodies rendered measurable by sensors, platforms, and companies; no checked source records the naming decision.
\end{itemize}
}

\section*{The Propers: Interpretive Possibilities}

\begin{studybox}{Editorial status}
The following exploratory proposals were developed in an AI-assisted editorial process. They are attributed to none of the cited authorities and claim no historical compositional intent. Each states its anchors and limits.
\end{studybox}

\Needspace{10\baselineskip}
\noindent\textbf{The whispering circuit is broken by rooftop speech.\enspace\elementref{First Reading; Acclamation; Gospel}}
Jeremiah hears hostile whispers and invitations to denounce; the acclamation names the Spirit's testimony and the disciples' derivative witness; Matthew moves concealed and darkened speech onto housetops. This spatial-verbal reversal breaks the persecutor's communications circuit: truth answers rumor not by imitation but by reception from God and public accountability. Separate reading misses the movement from enclosed surveillance through divine testimony to open acknowledgment. Its fruit for preaching, journalism, whistleblowing, and institutional truth-telling is courage joined to verifiability and nonretaliation. No checked source joined all three anchors. The Spirit's aid neither certifies every impulse, abolishes preparation or confidentiality, requires reckless disclosure, nor removes duties to protect sources and vulnerable persons.

\Needspace{10\baselineskip}
\noindent\textbf{Two kinds of watching assign opposite meanings to visibility.\enspace\elementref{First Reading; Gospel}}
Jeremiah's intimates watch for a misstep; Matthew's Father knows a sparrow's fall and a threatened disciple's hairs. Both minutely observe vulnerability, but hostile watching turns detail into leverage while divine knowing establishes worth and courage. Separate reading can leave Providence abstract and persecution external, missing the moral difference between surveillance awaiting failure and attention that makes no small life negligible. The fruit is an audit of pastoral records, digital monitoring, safeguarding data, and interpersonal attention: information serves care only where power answers to the person's good. Jerome and Chrysostom treat the strands separately; no exact precedent was located. Divine knowledge is not human data collection or positive causation of every fall and never excuses privacy, consent, or protection duties.

\Needspace{26\baselineskip}
\noindent\textbf{The one heard as poor is valued as sparrow and known as sheep.\enspace\elementref{Responsorial Psalm; Gospel; Communion B}}
Psalm 69 says God hears the needy; Matthew's slight sparrow denies that threatened witnesses are forgotten; Communion B, if appointed, names the Shepherd who knows his sheep and gives his life for them. The chain intensifies visibility into relation: hearing distress becomes recognition of irreducible worth, then self-giving knowledge. Separate reading misses divine attention becoming the pattern by which pastoral authority spends itself. The fruit is safeguarding and social care whose records, reporting, aid, and presence serve the vulnerable person rather than institutional reputation. Early shepherd texts offer near analogues, not this exact chain. Communion B is unresolved; Providence promises no immunity, “sheep” does not infantilize, and care requires competent systems rather than sentiment alone.

\Needspace{10\baselineskip}
\noindent\textbf{A surplus of grace can interrupt inherited fear.\enspace\elementref{Second Reading; Gospel}}
Romans traces sin and death through Adam yet makes Christ's gift more than an equal counterweight: grace abounds. Matthew names bodily threat, limits the killer's reach, and demands acknowledgment of Christ. Romans and Matthew are not officially correlated, yet this asymmetry suggests that inherited familial, cultural, or institutional fear is real yet not sovereign where a more abundant gift creates allegiance. Separate reading can leave Adam--Christ remote from public courage and “do not fear” to willpower; the connection grounds courage in received excess. Its fruit is communal resistance to fear's inherited scripts. Chrysostom's unequal-gift reading and ancient martyrial use of Matthew are near analogues, not this account. Paul offers no clinical or sociological theory; grace never replaces trauma-informed care, prudent flight, legal protection, or truthful danger assessment.

\Needspace{10\baselineskip}
\noindent\textbf{The Communion branches distinguish provision from recognition.\enspace\elementref{Communion A; Communion B}}
Communion A, if chosen, shows Psalm 145's eyes awaiting food; Communion B, if chosen, the Shepherd's knowledge and self-gift. Compared as alternatives rather than one enacted antiphon, they reveal care's two dimensions: dependable provision answers bodily dependence; personal knowledge and accountable presence resist reducing persons to need categories. Separate treatment notes Eucharistic aptitude but misses the branch structure's critique of systems without relationship and relationships without provision. Augustine on the psalm and early Shepherd reception offer near analogues, not this exact comparison. The fruit is integrated pastoral design. Selection is unknown, neither branch is a complete social program, and sacramental gift cannot be reduced to either humanitarian supply or therapeutic recognition.
